\documentclass[11pt]{article}

\usepackage[margin=2.5cm]{geometry}
\usepackage{graphicx}
\usepackage{xcolor}
\usepackage{booktabs}
\usepackage{caption}
\usepackage{enumitem}
\usepackage[hidelinks]{hyperref}

\graphicspath{{figures/}}

% Figure for a view. #1 description (unused), #2 image file stem in figures/,
% #3 width fraction, #4 caption.
\newcommand{\viewfig}[4]{%
  \begin{figure}[ht]\centering
  \includegraphics[width=#3\linewidth]{#2}
  \caption{#4}\label{fig:#2}
  \end{figure}%
}

\title{\textbf{Oceanus Folk Explorer}\\[2pt]
\large Visual Analysis of Musical Influence (VAST 2025, Mini-Challenge 1)}
\author{Iskandar Huseynzade\\\small Matricola 618017 \\ \small Visual Analytics (602AA)}
\date{}

\begin{document}
\maketitle

\begin{abstract}
\noindent Oceanus Folk Explorer is a web tool for exploring a large music knowledge graph.
The graph describes artists, songs, albums, labels, and the ways works influence one
another. The tool helps an analyst answer three questions: who shaped the career of the
star Sailor Shift, how the Oceanus Folk genre spread over time, and which artists are
likely to rise next. It uses three linked views built with D3.js: an influence network
placed on a time axis, a stacked-area chart of genre spread, and a trajectory chart with
a ranked list of rising stars. Selections and a time filter are shared across the views.
\end{abstract}

\section{Introduction}
The data comes from a fictional music scene on Oceanus Island. The central figure is
Sailor Shift, a singer who made the ``Oceanus Folk'' genre popular. The data is stored as
a knowledge graph: nodes are people, groups, songs, albums, and record labels, and edges
describe roles (who performed or wrote a work) and influence (which work drew from which).

An analyst working with this data has three goals:
\begin{enumerate}[noitemsep]
  \item Profile Sailor Shift: who influenced her over time, who she influenced, and her
        impact on the Oceanus Folk community.
  \item Understand the genre: how Oceanus Folk spread to other genres, whether the rise
        was steady or came in bursts, and which genres it drew from.
  \item Find rising stars: define what makes an artist promising, compare careers, and
        point to the next likely stars.
\end{enumerate}

The graph is large, so drawing all of it at once produces an unreadable tangle of lines,
often called a ``hairball''. The tool avoids this in two ways. It focuses on the part of
the graph around Sailor Shift, and it uses time as a strong layout rule so the picture
stays ordered. The interface is built with Vue 3, the charts with D3.js (version 7), and
the linking between views with crossfilter. Heavy computation is done once in Python and
exported as small JSON files that the browser loads quickly.

\section{The Data}
The graph is a directed multigraph with \textbf{17{,}412 nodes} and \textbf{37{,}857
edges}, split into 18 connected components. There are five node types and twelve edge
types. Table~\ref{tab:nodes} lists the node types and Table~\ref{tab:edges} the edges.

\begin{table}[ht]\centering
\caption{Node types and their main properties.}\label{tab:nodes}
\begin{tabular}{llr}
\toprule
Type & Main properties & Count \\
\midrule
Person & name, stage name & 11{,}361 \\
Song & name, genre, release year, notable, notoriety year & 3{,}615 \\
RecordLabel & name & 1{,}217 \\
Album & name, genre, release year, notable, notoriety year & 996 \\
MusicalGroup & name & 223 \\
\bottomrule
\end{tabular}
\end{table}

\begin{table}[ht]\centering
\caption{Edge types. Role edges connect people and works; influence edges connect a later
work to the earlier work it drew from.}\label{tab:edges}
\begin{tabular}{lll}
\toprule
Group & Edge type & Meaning (source $\rightarrow$ target) \\
\midrule
Role & PerformerOf & person/group performed a work \\
     & ComposerOf & person composed a work \\
     & ProducerOf & person/label produced a work or act \\
     & LyricistOf & person wrote lyrics for a work \\
     & MemberOf & person is a member of a group \\
     & RecordedBy & work was recorded by a label \\
     & DistributedBy & work was distributed by a label \\
\midrule
Influence & InStyleOf & work performed in the style of the target \\
          & InterpolatesFrom & work reuses a melody from the target \\
          & CoverOf & work is a cover of the target \\
          & LyricalReferenceTo & work refers to the target in its lyrics \\
          & DirectlySamples & work samples the target's audio \\
\bottomrule
\end{tabular}
\end{table}

\paragraph{Reading direction of influence.}
For the five influence edges, the \emph{target} is the earlier work that inspired, and the
\emph{source} is the later work that drew from it. In other words, the target is the
influencer. This single rule lets the tool place influence on a time axis, with older
works on the left and newer works on the right.

\paragraph{Time and popularity.}
Dates are stored as years. A work has a release year, and notable works also have a
``notoriety year'', which is the year the work first reached a top chart. The notable flag
and the notoriety year are the clearest signals of popularity in the data.

\paragraph{Derived data.}
Influence is recorded between works, not directly between people. To study people, the
tool builds person-to-person influence by following the creators of each work: if a work
by artist~A draws from a work by artist~B, then B influenced A. Three small files are
produced in Python:
\begin{itemize}[noitemsep]
  \item \texttt{sailor\_ego.json} -- Sailor Shift's influence neighborhood (138 artists and
        541 influence links at one hop), with a year and simple metrics for each artist.
  \item \texttt{genre\_spread.json} -- yearly counts of works influenced by Oceanus Folk,
        grouped by genre, plus the genres it drew from and the artists it affected.
  \item \texttt{rising\_stars.json} -- a Rising Star Score and a career trajectory for each
        Oceanus Folk artist.
\end{itemize}

The data has a few rough edges that the preprocessing handles. Some influence edges point
at a person or a group instead of a work, so both ends of every influence edge are
resolved to people before counting. A small number of works have no year and are left out
of time-based charts.

\section{Related Work}
Two classic ways to show a graph are the node-link diagram and the adjacency matrix.
Node-link diagrams are easy to read when small but become tangled when large. Matrices
stay readable when dense but hide paths. Hybrid designs such as NodeTrix~\cite{nodetrix}
mix both. To keep the focus clear, many tools show one node and its neighborhood, which is
the approach used here.

Graphs that change over time need an extra cue for time. A survey by Beck et
al.~\cite{dyngraph} groups these methods into animation and timeline-based layouts. Placing
time on an axis, as this tool does for influence, is a timeline-based choice that keeps the
order of events visible without animation.

To show how parts of a whole change over the years, the stacked graph, or
streamgraph~\cite{stackedgraphs}, is a common choice and has been used for music listening
histories. The genre-spread view uses a stacked area of the same family. For working
across several views at once, coordinated multiple views with linked selection and
brushing~\cite{cmv} are a standard pattern, and the tool follows it. All charts are drawn
with D3~\cite{d3}.

\section{Design Choices}
\paragraph{Scope.}
The tool starts from Sailor Shift and shows her neighborhood, not the whole graph. This
keeps the first screen readable and matches the main analytical question.

\paragraph{Layout.}
In the influence network, the horizontal position of an artist is the year they became
active. A force simulation only spreads nodes vertically and keeps them from overlapping,
while a strong horizontal force pins each node to its year. The result reads from left to
right as a flow of influence through time, instead of a shapeless cloud.

\paragraph{Color.}
Genre is the main color channel. Oceanus Folk uses a clear cyan so it stands out, and other
genres use a categorical palette. Influence edges are colored by their type. In the rising
stars view, the five parts of the score each have a fixed color so the same color always
means the same thing.

\paragraph{Size and shape.}
Node size in the network grows with the number of other artists an artist influenced, so
important artists are larger. Bars and areas use length and height, which are easy to
compare.

\paragraph{Transformations.}
Career trajectories are cumulative sums, so a steeper line means faster growth. Influence
reach is passed through a logarithm so that one very prolific artist does not flatten the
scale for everyone else. Score parts are normalized to the range zero to one before they
are combined.

\paragraph{Readability.}
Only the focus artist is labeled in the network. Full details appear in a side panel when a
node is selected, which keeps the chart clean while still giving access to the numbers.

\section{The Tool}
The tool has three views, chosen with a tab at the top. The views share the current
selection and the current year filter.

\subsection{Artist View}
This view answers the first question. The center is the influence network described above.
\viewfig{artist view}{artist}{0.9}{The artist view: influence network on a time axis, a
year timeline, a ranking of top influencers, and a detail panel.}
Older influencers sit to the left of Sailor Shift and the artists she influenced sit to the
right. A timeline below the network shows how many artists were active each year and can be
brushed to filter the other views by year. A toggle in the header switches between the
one-hop and two-hop neighborhood so the analyst can broaden the view at the cost of density. A ranking lists the artists who influenced the
most others. A detail panel shows the selected artist's genre, active years, number of
works, and the lists of who influenced them and who they influenced. Clicking a node, a
ranking row, or a name in the panel selects that artist and updates every part of the view.

\subsection{Genre Spread View}
This view answers the second question.
\viewfig{genre spread view}{genre}{0.9}{The genre spread view: a stacked area of works
influenced by Oceanus Folk over time, with rankings of genres and affected artists.}
The main chart is a stacked area. For each year it shows how many works that were influenced
by Oceanus Folk appeared, split by the genre of those works. A dashed line shows Oceanus
Folk's own releases for comparison. The shape of the area answers whether the rise was
steady or came in bursts: the data shows clear peaks around 2017, 2020, and a large one in
2023, so the rise was intermittent. Side panels rank the genres Oceanus Folk drew from
(Indie Folk is by far the largest, followed by Synthwave and Dream Pop), the genres it
influenced (Desert Rock, Indie Folk, Dream Pop, and Space Rock lead), and the artists most
affected by it.

\subsection{Rising Stars View}
This view answers the third question.
\viewfig{rising stars view}{stars}{0.9}{The rising stars view: career trajectories on the
left and a ranked leaderboard with a score breakdown on the right.}
The left chart compares career trajectories. Each line is an artist's cumulative number of
chart hits over time, so a steep recent slope marks a fast riser. Sailor Shift is drawn as
a dashed benchmark line. The right side is a leaderboard of candidates, ranked by a Rising
Star Score. The score is the weighted sum of five parts: recent chart hits, momentum
(recent hits minus earlier hits), influence reach, share of Oceanus Folk works, and career
youth. Each row's bar is split into these five parts in their fixed colors, so the analyst
can see why an artist ranks where it does. Clicking a candidate adds or removes its line
from the trajectory chart.

\subsection{Methodology: the Rising Star Score}
The score combines five normalized parts. All parts are in the range $[0,1]$ and the final
score is the weighted sum. The defaults in Table~\ref{tab:weights} are loaded from the data
file, and sliders in the right panel let the analyst change any weight live; the leaderboard
re-ranks at once. This makes the prediction transparent and lets the analyst test how
sensitive the ranking is to each criterion.

\begin{table}[ht]\centering
\caption{Components and default weights of the Rising Star Score.}\label{tab:weights}
\begin{tabular}{lll}
\toprule
Part & Meaning & Weight \\
\midrule
Recent chart hits & notable works in the last 5 years (popularity) & 0.30 \\
Momentum & recent hits minus earlier hits (acceleration) & 0.20 \\
Influence reach & distinct downstream artists, log-damped & 0.20 \\
Oceanus Folk affinity & share of the artist's works in the genre & 0.20 \\
Career youth & shorter career = more upside & 0.10 \\
\bottomrule
\end{tabular}
\end{table}

Influence reach is passed through $\log(1+x)$ before normalization so that one outlier
artist credited on hundreds of works does not flatten the scale for the others. To be a
candidate at all an artist must have at least one Oceanus Folk work, a known debut year,
and either a notable work or a reach of two or more.

\subsection{Linking}
The views are connected in three ways. Selecting an artist in the artist view updates the
network highlight, the ranking, and the detail panel together. Brushing a year range in the
timeline filters the network and the ranking through crossfilter. Selection is also shared
across tabs: a rising-star candidate that exists in Sailor Shift's neighborhood has a small
arrow button that opens the artist view focused on that artist, and the leaderboard
highlights whichever artist is currently selected in the network. The tab bar switches the
analytical task while keeping the same underlying data in memory.

\section{Use Case}
Suppose an analyst wants to trace Sailor Shift's career and then look ahead. They open the
artist view. Sailor Shift sits near 2024. To her left are the artists who influenced her,
such as Coralia Bellweather and Barnaby Quill, connected by ``in the style of'' and
``interpolates from'' links. To her right are the artists she influenced, such as Bea Azure
and Zara Quinn. Selecting Coralia Bellweather shows that she is both an early influence on
Sailor and a strong influencer in her own right.

The analyst then switches to the genre spread view. The stacked area shows that Oceanus
Folk did not grow smoothly; it grew in bursts, with a large jump in 2023. The side panels
show that the genre drew most heavily from Indie Folk and in turn spread into Desert Rock,
Dream Pop, and Space Rock. This explains how the genre evolved by borrowing from and
feeding back into nearby styles.

Finally, the analyst opens the rising stars view. Sailor Shift's trajectory is the highest,
which fits her role as the established star. The leaderboard ranks other Oceanus Folk
artists by the Rising Star Score, and the breakdown bars show that some are ranked for their
recent chart hits while newer artists such as Arlo Sterling rank for youth and momentum.
Adding a few of these artists to the trajectory chart lets the analyst compare their growth
directly and pick the most promising ones.

\section{Conclusion}
Oceanus Folk Explorer turns a large and tangled knowledge graph into three readable views
that answer concrete questions about influence, genre spread, and rising stars. The main
design idea is to use time as the backbone of the layout and to focus on the area around
Sailor Shift, which keeps the graph clear.

The tool has limits. The influence edges are sparse and were collected by hand, so some
relationships are missing. Works with many creators can inflate influence counts, which the
score reduces with a logarithm but does not fully remove. The data also includes future
years, so predictions should be read as a guide rather than a fact. Useful next steps are a
view for the collaboration network (the data is already prepared but not yet drawn) and a
way to compare two artists side by side beyond the trajectory chart.

\begin{thebibliography}{9}
\bibitem{nodetrix} N.~Henry, J.-D.~Fekete, and M.~J.~McGuffin. NodeTrix: A Hybrid
Visualization of Social Networks. \emph{IEEE Transactions on Visualization and Computer
Graphics}, 2007.
\bibitem{dyngraph} F.~Beck, M.~Burch, S.~Diehl, and D.~Weiskopf. A Taxonomy and Survey of
Dynamic Graph Visualization. \emph{Computer Graphics Forum}, 2017.
\bibitem{stackedgraphs} L.~Byron and M.~Wattenberg. Stacked Graphs -- Geometry and
Aesthetics. \emph{IEEE Transactions on Visualization and Computer Graphics}, 2008.
\bibitem{cmv} J.~C.~Roberts. State of the Art: Coordinated and Multiple Views in
Exploratory Visualization. \emph{International Conference on Coordinated and Multiple Views
in Exploratory Visualization}, 2007.
\bibitem{d3} M.~Bostock, V.~Ogievetsky, and J.~Heer. D3: Data-Driven Documents.
\emph{IEEE Transactions on Visualization and Computer Graphics}, 2011.
\end{thebibliography}

\end{document}
